\section{INTRODUCCIÓN}
Los procesos electorales constituyen uno de los mecanismos mediante los cuales las sociedades democráticas expresan su voluntad colectiva, y su legitimidad depende de la integridad y la confidencialidad de la información que generan (Baloglu et al., 2021).

La seguridad de la votación electrónica se ha estudiado principalmente sobre esquemas de clave pública como RSA. Baloglu et al. (2021) emplearon verificación formal para analizar propiedades de integridad, confidencialidad y autenticidad en Helios y Belenios, evidenciando que estas dependen tanto de las primitivas criptográficas como del diseño del protocolo. Sin embargo, estos esquemas son vulnerables, en teoría, al algoritmo cuántico de Shor. Como respuesta, el NIST publicó en 2024 los estándares FIPS 203 y FIPS 204, sustentados en problemas sobre retículos algebraicos --- Module-LWE y Module-SIS.

Esta transición hacia primitivas post-cuánticas aún no se refleja en los antecedentes sobre votación electrónica basada en blockchain, incluidos los del ámbito universitario boliviano (Churata Sonco, 2020; Fernández Tristán, 2021), lo que evidencia una brecha vigente incluso a nivel internacional.

Ramos y Gómez (2023) realizaron una revisión sistemática de protocolos de votación electrónica basados en criptografía post-cuántica, sistematizando las propuestas reportadas en la literatura internacional. Dicho trabajo se circunscribe a la síntesis del estado del arte y no propone, diseña ni verifica formalmente un protocolo propio, ni aborda el contexto universitario boliviano; constituye un antecedente secundario que evidencia el interés creciente por el área, sin cubrir el vacío relativo al diseño y verificación formal de un protocolo específico en dicho contexto.

El presente trabajo, \textit{Protocolo para la votación electrónica basado en criptografía post-cuántica de retículos}, se sitúa en dicho espacio, estructurado bajo \textit{Design Science Research} de Vaishnavi y Kuechler (2004): se diseña un protocolo con las primitivas ML-KEM y ML-DSA, modelado en \textit{Tamarin Prover} bajo el modelo Dolev--Yao, y se contrasta cuasi-experimentalmente con un protocolo equivalente basado en RSA/ECDSA.
